\section{从开环到闭环反馈}

一台恒温水箱要把水温稳定在 \(60\) 度。设水温为 \(x(t)\)，环境温度 \(x_e=20\) 度，散热使水温按 \(-a(x-x_e)\) 的速率向环境滑落（\(a>0\)），加热器的贡献记为输入 \(u\)：

\[
\dot x=-a\bigl(x-x_e\bigr)+u.
\]

怎样选 \(u\)，让水温到达并停在 \(60\) 度？本节用这个例子对比两种答案，并回答一个更尖锐的问题：为什么``把计划一次算好再执行''在现实中几乎必然失败，而``边看边调''却稳得住。

\subsection{开环：事先算好全部输入}

设 \(a=0.1\)。要让水温停在目标 \(r=60\)，稳态要求 \(0=-0.1(60-20)+u\)，即保持功率 \(u^\ast=4\)；再配一段升温用的输入曲线，整个控制可以在机器开动之前一次算完。这就是开环控制：输入是时间的函数 \(u(t)\)，执行期间不再看水温。

记偏差 \(e=x-r\)，把输入拆成``计划保持量''与``修正量''：\(u=u^\ast+\bar u\)。代入方程后常数项恰好相消，偏差满足

\[
\dot e=-a\,e+\bar u.
\]

开环方案取 \(\bar u\equiv0\)。若模型精确、世界安静，则 \(e\to0\)，一切完美。

\subsection{认知冲突：计划赶不上现实}

现实从两个方向破坏这个完美。

第一，模型不可能精确。真实散热系数若是 \(\tilde a=0.08\) 而不是 \(0.1\)，同一个保持功率对应的真实平衡点变成 \(x=70\) 度——计划中的 \(60\) 度在真实系统里根本不是平衡点，而执行者对此一无所知，因为它从不测量。

第二，扰动不请自来。开了一扇窗，方程变成

\[
\dot x=-a\bigl(x-x_e\bigr)+u+w,\qquad w=-1,
\]

偏差方程变为 \(\dot e=-ae+w\)，稳态偏差 \(e=w/a=-10\) 度：水温停在 \(50\) 度，计划依然``以为''一切正常。

更危险的是被控对象本身带正反馈的情形（例如放热反应：温度越高反应越快），偏差方程形如 \(\dot e=+\lambda e+w\)（\(\lambda>0\)）：计划与现实的任何微小差异都被指数放大。``离散动力系统、分岔与混沌'' 中 Logistic 映射对初值的敏感性，是同一现象的离散版本——那一篇的教训是：指数发散面前，任何事先算好的轨线都只有有限的有效期。

开环的根本缺陷不是算得不准，而是结构上封死了纠正的可能：输入里不含任何关于现实的信息。

\subsection{闭环：让误差自己产生纠偏力}

闭环控制只做一处改动：把测量到的偏差接回输入，

\[
\bar u=-k\,e,\qquad k>0,
\]

即水温低于目标就加大功率，高于目标就减小。偏差方程变为

\[
\dot e=-(a+k)\,e+w.
\]

两件事立刻改变。其一，在模型准确且无扰动时 \(e\to0\)，而且收敛率从 \(a\) 提高到 \(a+k\)。其二，常值扰动下的稳态偏差从 \(w/a\) 缩小到

\[
e_{\mathrm{ss}}=\frac{w}{a+k},
\]

增益 \(k\) 越大，残余偏差越小。开窗的例子代入 \(k=1\)：\(e_{\mathrm{ss}}=-1/1.1\approx-0.9\) 度，对比开环的 \(-10\) 度。模型失配也会在偏差方程中表现成等效扰动，所以比例反馈同样只能压小其影响，不能保证精确归零；下文的积分项才负责消除这类常值残差。

为什么``接回误差''有这样本质的效果？用 ``稳定性与 Lyapunov 函数'' 的语言最清楚。取 \(V(e)=\tfrac12 e^2\)，它在目标处取最小、处处非负；无扰动时沿闭环轨迹

\[
\dot V=e\,\dot e=-(a+k)e^2<0
\qquad(e\ne0).
\]

\(V\) 扮演``高度''的角色并被动力学严格消耗，偏差只能向零沉——反馈把闭环系统改造成了一个以目标为汇的梯度流。开环没有这个性质：它的 \(V\) 只被扰动推着走，没有任何力把它拉回来。

\subsection{PID：三种纠偏力的分工}

工程界最常用的反馈律把纠偏力拆成三项：

\[
u(t)=u^\ast-k_P\,e(t)-k_I\int_0^t e(\tau)\,d\tau-k_D\,\dot e(t),
\]

比例（P）、积分（I）、微分（D）各管一段时间尺度，各治一种病，也各带一种副作用。

\textbf{比例管现在。}偏差多大就纠多大，直观、见效快。但它的病上面已经看到：有常值扰动时，稳态必须留下 \(w/(a+k)\) 的偏差来``养活''纠偏力——偏差归零的瞬间比例项也为零，功率随之中断，扰动又把系统推开。比例控制无法彻底消除静差；而一味加大 \(k_P\) 会吃掉稳定裕度，把迟缓的纠正变成剧烈的振荡。

\textbf{积分管历史。}只要偏差没有归零，积分项就持续累积，把输入一路推高，直到偏差被彻底抹平——它治的正是比例留下的静差。代价是新的病：执行器有功率上限，饱和期间误差继续累积，积分项胀得过大，等误差反号后还要花很长时间``退账''，造成大幅超调，这就是积分饱和（integrator windup），工程上要对积分项限幅（anti-windup）。积分还引入额外的相位滞后，参数不当时本身就能让系统振荡。

\textbf{微分管趋势。}它不看偏差有多大，只看偏差变化有多快：水温正快速冲向目标时提前收力，相当于给系统加阻尼，抑制超调。它的病在于噪声：测量噪声多为高频抖动，微分恰好放大高频，\(k_D\) 一大，控制信号就被噪声淹没，实际使用中微分项必须配低通滤波。

三项的分工因此是：P 给即时纠偏，I 消灭残差，D 预判趋势。调参的本质是在``纠得快''与``稳得住、抗噪''之间权衡——PID 里没有``最优''的概念，只有够用的权衡。

\subsection{反馈不是万能的}

反馈只能对``已经发生的偏差''行动：测量有延迟，纠偏就有延迟，延迟配上过大的增益会把稳定系统逼成振荡。反馈还把测量噪声直接写进控制信号。更根本地，反馈默认``目标已知、偏差可测''——目标怎么选、多个目标之间怎样权衡，它一概不管。实践中常见的做法是开环与闭环混合：用模型算一个大致正确的前馈（feedforward）计划，再用反馈清理残差。

承认这些边界后，还剩一个 PID 回答不了的问题：什么叫``好''的控制？稳住只是及格线；省燃料、省时间、少超调，这些要求必须先写成目标函数，才谈得上优化。这是下一节的主题。比例控制需要非零偏差来维持抵消常值扰动的控制量，所以必留静差；积分项能对常值失配累积出固定补偿，却不能自动精确跟踪任意时变扰动，而且还会引入相位滞后。这正说明控制结构必须与扰动类型匹配。
